\Chapter{42}

\Verse{1}
{
    \zA{话说贾母王夫人去后，姐妹们复进园来吃饭。}
}
{
    \eA{We will now resume our story by adding that, on the return of the young
        ladies into the garden, they had their meal. This over,}
}
{
}

\Verse{2}
{
    \zA{那刘老老带着板儿，先来见凤姐 儿说：“明日一早定要家去了。}
}
{
    \eA{they parted company, and nothing more need be said about them. We will
        notice, however, that old goody Liu took Pan Erh along with her, and came
        first and paid a visit to lady Feng.}
}
{
}

\Verse{3}
{
    \zA{虽然住了两三天，日子却不多，把古往今来没见过 的、没吃过的、没听见的都经验过了。}
}
{
    \eA{"We must certainly start for home to-morrow, as soon as it is daylight," she
        said. "I've stayed here, it's true, only two or three days, but in these few
        days I have reaped experience in everything that I had not seen from old
        till now.}
}
{
}

\Verse{4}
{
    \zA{难得老太太和姑奶奶并那些小姐们，连各房 里的姑娘们，都这样怜贫惜老照看我。}
}
{
    \eA{It would be difficult to find any one as compassionate of the poor and
        considerate to the old as your venerable dame, your Madame Wang, your young
        ladies, and the girls too attached to the various rooms, have all shown
        themselves in their treatment of me!}
}
{
}

\Verse{5}
{
    \zA{我这一回去没别的报答，惟有请些高香，天 天给你们念佛，保佑你们长命百岁的，就算我的心了。”}
}
{
    \eA{When I get home now, I shall have no other means of showing how grateful I
        am to you than by purchasing a lot of huge joss-sticks and saying daily
        prayers to Buddha on your behalf; and if he spares you all to enjoy a long
        life of a hundred years my wishes will be accomplished."}
}
{
}

\Verse{6}
{
    \zA{凤姐儿笑道：“你别喜欢， 都是为你，老太太也叫风吹病了，躺着嚷不舒服；我们大姐儿也着了凉了，在那里 发热呢。”}
}
{
    \eA{"Don't be so exultant!" lady Feng smilingly replied. "It's all on account of
        you that our old ancestor has fallen ill, by exposing herself to draughts
        and that she suffers from disturbed sleep; also that our Ta Chieh-erh has
        caught a chill and is laid up at home with fever." Goody Liu, at these
        words,}
}
{
}

\Verse{7}
{
    \zA{刘老老听了，忙叹道：“老太太有年纪了，不惯十分劳乏的。”}
}
{
    \eA{speedily heaved a sigh. "Her venerable ladyship," she said, "is a person
        advanced in years and not accustomed to any intense fatigue!" "She has never
        before been in such high spirits as yesterday!"}
}
{
}

\Verse{8}
{
    \zA{凤姐儿道： “从来不像昨儿高兴。}
}
{
    \eA{lady Feng observed. "As you were here, so anxious was she to let you see
        everything,}
}
{
}

\Verse{9}
{
    \zA{往常也进园子逛去，不过到一两处坐坐就来了。}
}
{
    \eA{that she trudged over the greater part of the garden. And Ta Chieh-erh was
        given a piece of cake by Madame Wang, when I came to hunt you up,}
}
{
}

\Verse{10}
{
    \zA{昨儿因为你 在这里，要叫都逛逛，一个园子倒走了多半个。}
}
{
    \eA{and she ate it, who knows in what windy place, and began at once to get
        feverish." "Ta Chieh-erh," goody Liu remarked, "hasn't, I fancy, often put
        her foot into the garden;}
}
{
}

\Verse{11}
{
    \zA{大姐儿因为我找你去，太太递了一 块糕给他，谁知风地里吃了，就发起热来。”}
}
{
    \eA{and young people like her mustn't really go into strange places, for she's
        not like our children, who are able to use their legs! In what graveyards
        don't they ramble about! A puff of wind may, on the one hand,}
}
{
}

\Verse{12}
{
    \zA{刘老老道：“妞妞儿只怕不大进园子。}
}
{
    \eA{have struck her, it's not at all unlikely; or being, on the other, so chaste
        in body, and her eyes also so pure she may,}
}
{
}

\Verse{13}
{
    \zA{比不得我们的孩子，一会走，那个坟圈子里不跑去？}
}
{
    \eA{it is to be feared, have come across some spirit or other. I can't help
        thinking therefore that you should consult some book of exorcisms on her
        behalf;}
}
{
}

\Verse{14}
{
    \zA{一则风拍了也是有的，二则只 怕他身上干净，眼睛又净，或是遇见什么神了。}
}
{
    \eA{for mind she may have run up against some evil influence." This remark
        suggested the idea to lady Feng. There and then she called P'ing Erh to
        fetch the 'Jade Box Record.' When brought, she desired Ts'ai Ming to look
        over it for her.}
}
{
}

\Verse{15}
{
    \zA{依我说，给他瞧瞧祟书本子，仔细 撞客着。”}
}
{
    \eA{Ts'ai Ming turned over the pages for a time, and then read: 'Those who fall
        ill on the 25th day of the 8th moon have come across,}
}
{
}

\Verse{16}
{
    \zA{一语提醒了凤姐儿，便叫平儿拿出《玉匣记》来，叫彩明来念。}
}
{
    \eA{in a due westerly quarter, of some flower spirit; they feel heavy, with no
        inclination for drink or food. Take seven sheets of white paper money,}
}
{
}

\Verse{17}
{
    \zA{彩明翻 了一会子，念道：“八月二十五日病者，东南方得之，有缢死家亲女鬼作祟，又遇 花神。}
}
{
    \eA{and, advancing forty steps due west, burn them and exorcise the spirit;
        recovery will follow at once!'" "There's really no mistake about that!" lady
        Feng smiled. "Are there not flower spirits in the garden?}
}
{
}

\Verse{18}
{
    \zA{用五色纸钱四十张，向东南方四十步送之大吉。”}
}
{
    \eA{But what I dread is that our old lady mayn't have come across one too."
        Saying this, she bade a servant purchase two lots of paper money.}
}
{
}

\Verse{19}
{
    \zA{凤姐儿笑道：“果然不错， 园子里头可不是花神!只怕老太太也是遇见了。”}
}
{
    \eA{On their arrival, she sent for two proper persons, the one to exorcise the
        spirits for dowager lady Chia and the other to expel them from Ta Chieh-erh;
        and these observances over, Ta Chieh-erh did, in effect,}
}
{
}

\Verse{20}
{
    \zA{一面命人请两分纸钱来，着两个人 来，一个与贾母送祟，一个与大姐儿送祟，果见大姐儿安稳睡了。}
}
{
    \eA{drop quietly to sleep. "It's verily people advanced in years like you," lady
        Feng smilingly exclaimed; "who've gone through many experiences! This Ta
        Chieh-erh of mine has often been inclined to ail, and it has quite puzzled
        me to make out how and why it was."}
}
{
}

\Verse{21}
{
    \zA{凤姐儿笑道：“到底是你们有年纪的经历的多。}
}
{
    \eA{"This isn't anything out of the way!" goody Liu said. "Affluent and
        honourable people bring up their offspring to be delicate.}
}
{
}

\Verse{22}
{
    \zA{我们大姐儿时常肯病，也不知 是什么原故。”}
}
{
    \eA{So naturally, they are not able to endure the least hardship! Moreover, that
        young child of yours is so excessively cuddled that she can't stand it.}
}
{
}

\Verse{23}
{
    \zA{刘老老道：“这也有的。}
}
{
    \eA{Were you, therefore, my lady, to pamper her less from henceforth,}
}
{
}

\Verse{24}
{
    \zA{富贵人家养的孩子都娇嫩，自然禁不得一些 儿委屈。}
}
{
    \eA{she'll steadily improve." "There's plenty of reason in that too!" lady Feng
        observed. "But it strikes me that she hasn't as yet got a name,}
}
{
}

\Verse{25}
{
    \zA{再他小人儿家，过于尊贵了也禁不起。}
}
{
    \eA{so do give her one in order that she may borrow your long life! In the next
        place, you are country-people,}
}
{
}

\Verse{26}
{
    \zA{以后姑奶奶倒少疼他些就好了。”}
}
{
    \eA{and are, after all,—I don't expect you'll get angry when I mention
        it,—somewhat in poor circumstances.}
}
{
}

\Verse{27}
{
    \zA{凤姐儿道：“也是有的。}
}
{
    \eA{Were a person then as poor as you are to suggest a name for her,}
}
{
}

\Verse{28}
{
    \zA{我想起来，他还没个名字，你就给他起个名字，借借你的 寿；二则你们是庄家人，不怕你恼，到底贫苦些，你们贫苦人起个名字只怕压的住。”}
}
{
    \eA{you may, I trust, have the effect of counteracting this influence for her."
        When old goody Liu heard this proposal, she immediately gave herself up to
        reflection. "I've no idea of the date of her birth!" she smiled after a
        time. "She really was born on no propitious date!" lady Feng replied. "By a
        remarkable coincidence she came into the world on the seventh day of the
        seventh moon!"}
}
{
}

\Verse{29}
{
    \zA{刘老老听说，便想了一想，笑道：“不知他是几时养的？”}
}
{
    \eA{"This is certainly splendid!" old goody Lin laughed with alacrity. "You had
        better name her at once Ch'iao Chieh-erh (seventh moon and ingenuity).}
}
{
}

\Verse{30}
{
    \zA{凤姐儿道：“正是养的日 子不好呢：可巧是七月初七日。”}
}
{
    \eA{This is what's generally called: combating poison by poison and attacking
        fire by fire. If therefore your ladyship fixes upon this name of mine,}
}
{
}

\Verse{31}
{
    \zA{刘老老忙笑道：“这个正好，就叫做巧姐儿好。}
}
{
    \eA{she will, for a surety, attain a long life of a hundred years; and when she
        by and bye grows up to be a big girl, every one of you will be able to have
        a home and get a patrimony!}
}
{
}

\Verse{32}
{
    \zA{这 个叫做‘以毒攻毒，以火攻火’的法子。}
}
{
    \eA{Or if, at any time, there occur anything inauspicious and she has to face
        adversity, why it will inevitably change into prosperity;}
}
{
}

\Verse{33}
{
    \zA{姑奶奶定依我这名字，必然长命百岁。}
}
{
    \eA{and if she comes across any evil fortune, it will turn into good fortune.
        And this will all arise from this one word,}
}
{
}

\Verse{34}
{
    \zA{日 后大了，各人成家立业，或一时有不遂心的事，必然遇难成祥，逢凶化吉，都从这 ‘巧’字儿来。”}
}
{
    \eA{'Ch'iao' (ingenuity.)" Lady Feng was, needless to say, delighted by what she
        heard, and she lost no time in expressing her gratitude. "If she be
        preserved," she exclaimed, "to accomplish your good wishes,}
}
{
}

\Verse{35}
{
    \zA{凤姐儿听了，自是欢喜，忙谢道：“只保佑他应了你的话就好了。”}
}
{
    \eA{it will be such a good thing!" Saying this, she called P'ing Erh. "As you
        and I are bound to be busy to-morrow," she said, "and won't, I fear, be able
        to spare any leisure moments,}
}
{
}

\Verse{36}
{
    \zA{说着，叫平儿来吩咐道：“明儿咱们有事，恐怕不得闲儿，你这会子闲着，把送老 老的东西打点了，他明儿一早就好走的便宜了。”}
}
{
    \eA{you'd better, if you have nothing to do now, get ready the presents for old
        goody Liu, so as to enable her to conveniently start at early dawn
        to-morrow." "How could I presume to be the cause of such reckless waste?"
        goody Liu interposed. "I've already disturbed your peace and quiet for
        several days,}
}
{
}

\Verse{37}
{
    \zA{刘老老道：“不敢多破费了。}
}
{
    \eA{and were I to also take your things away, I'd feel still less at ease in my
        heart!"}
}
{
}

\Verse{38}
{
    \zA{已经 遭扰了几天，又拿着走，越发心里不安了。”}
}
{
    \eA{"There's nothing much!" lady Feng protested. "They consist simply of a few
        ordinary things. But, whether good or bad,}
}
{
}

\Verse{39}
{
    \zA{凤姐儿笑道：“也没有什么，不过随常 的东西。}
}
{
    \eA{do take them along, so that the people in the same street as yourselves and
        your next-door neighbours may have some little excitement,}
}
{
}

\Verse{40}
{
    \zA{好也罢，歹也罢，带了去，你们街坊邻舍看着也热闹些，也是上城一趟。”}
}
{
    \eA{and that it may look as if you had been on a visit to the city!" But while
        she endeavoured to induce the old dame to accept the presents, she noticed
        P'ing Erh approach. "Goody Liu," she remarked,}
}
{
}

\Verse{41}
{
    \zA{说着只见平儿走来说：“老老过这边瞧瞧。”}
}
{
    \eA{"come over here and see!" Old goody Liu precipitately followed P'ing Erh
        into the room on the off side.}
}
{
}

\Verse{42}
{
    \zA{刘老老忙跟了平儿到那边屋里，只见堆 着半炕东西。}
}
{
    \eA{Here she saw the stove-couch half full with piles of things. P'ing Erh took
        these up one by one and let her have a look at them.}
}
{
}

\Verse{43}
{
    \zA{平儿一一的拿给他瞧着，又说道：“这是昨日你要的青纱一匹，奶奶 另外送你一个实地月白纱做里子。}
}
{
    \eA{"This," she explained, "is a roll of that green gauze you asked for
        yesterday. Besides this, our lady Feng gives you a piece of thick
        bluish-white gauze to use as lining. These are two pieces of pongee, which
        will do for wadded coats and jupes as well. In this bundle are two pieces of
        silk,}
}
{
}

\Verse{44}
{
    \zA{这是两个茧绸，做袄儿裙子都好。}
}
{
    \eA{for you to make clothes with, for the end of the year. This is a box
        containing various home-made cakes.}
}
{
}

\Verse{45}
{
    \zA{这包袱里是两 匹绸子，年下做件衣裳穿。}
}
{
    \eA{Among them are some you've already tasted and some you haven't; so take them
        along, and put them in plates and invite your friends;}
}
{
}

\Verse{46}
{
    \zA{这是一盒子各样内造小饽饽儿，也有你吃过的，也有没 吃过的，拿去摆碟子请人，比买的强些。}
}
{
    \eA{they'll be ever so much better than any that you could buy! These two bags
        are those in which the melons and fruit were packed up yesterday. This one
        has been filled with two bushels of fine rice, grown in the imperial fields,
        the like of which for congee,}
}
{
}

\Verse{47}
{
    \zA{这两条口袋是你昨日装果子的，如今这一 个里头装了两斗御田粳米，熬粥是难得的；这一条里头是园子里的果子和各样干果 子。}
}
{
    \eA{it would not be easy to get. This one contains fruits from our garden and
        all kinds of dry fruits. In this packet, you'll find eight taels of silver.
        These various things are presents for you from our Mistress Secunda. Each of
        these packets contains fifty taels so that there are in all a hundred taels;
        they're the gift of Madame Wang.}
}
{
}

\Verse{48}
{
    \zA{这一包是八两银子。}
}
{
    \eA{She bids you accept them so as to either carry on any trade,}
}
{
}

\Verse{49}
{
    \zA{这都是我们奶奶的。}
}
{
    \eA{for which no big capital is required, or to purchase several acres of land,}
}
{
}

\Verse{50}
{
    \zA{这两包每包五十两，共是一百两，是 太太给的，叫你拿去，或者做个小本买卖，或者置几亩地，以后再别求亲靠友的。”}
}
{
    \eA{in order that you mayn't henceforward have any more to beg favours of
        relatives, or to depend upon friends." Continuing, she added smilingly, in a
        low tone of voice, "These two jackets, two jupes, four head bands, and a
        bundle of velvet and thread are what I give you, worthy dame, as my share.}
}
{
}

\Verse{51}
{
    \zA{说着又悄悄笑道：“这两件袄儿和两条裙子，还有四块包头，一包绒线，可是我送 老老的。}
}
{
    \eA{These clothes are, it is true, the worse for use, yet I haven't worn them
        very much. But if you disdain them, I won't be so presuming as to say
        anything." After mention of each article by P'ing Erh, goody Liu muttered
        the name of Buddha,}
}
{
}

\Verse{52}
{
    \zA{那衣裳虽是旧的，我也没大很穿，你要弃嫌，我就不敢说了。”}
}
{
    \eA{so already she had repeated Buddha's name several thousands of times. But
        when she saw the heap of presents which P'ing Erh too bestowed on her, and
        the little ostentation with which she did it, she promptly smiled. "Miss!"}
}
{
}

\Verse{53}
{
    \zA{平儿说一样，刘老老就念一句佛，已经念了几千佛了；又见平儿也送他这些东 西，又如此谦逊，忙笑道：“姑娘说那里话？}
}
{
    \eA{she said, "what are you saying? Could I ever disdain such nice gifts as
        these! Had I even the money, I couldn't buy them anywhere. The only thing is
        that I feel overpowered with shame. If I keep them, it won't be nice, and if
        I don't accept them, I shall be showing myself ungrateful for your kind
        attention."}
}
{
}

\Verse{54}
{
    \zA{这样好东西，我还弃嫌!我就有银子， 没处买这样的去呢。}
}
{
    \eA{"Don't utter all this irrelevant talk!" P'ing Erh laughed. "You and I are
        friends; so compose your mind and take the things I gave you just now!
        Besides, I have, on my part,}
}
{
}

\Verse{55}
{
    \zA{只是我怪臊的，收了不好，不收又辜负了姑娘的心。”}
}
{
    \eA{something to ask of you. When the close of the year comes, select a few of
        your cabbages, dipped in lime, and dried in the sun,}
}
{
}

\Verse{56}
{
    \zA{平儿笑 道：“别说外话，咱们都是自己，我才这么着。}
}
{
    \eA{as well as some lentils, flat beans, tomatoes and pumpkin strips, and
        various sorts of dry vegetables and bring them over. We're all, both high or
        low,}
}
{
}

\Verse{57}
{
    \zA{你放心收了罢，我还和你要东西呢。}
}
{
    \eA{fond of such things. These will be quite enough! We don't want anything
        else, so don't go to any useless trouble!"}
}
{
}

\Verse{58}
{
    \zA{到年下，你只把你们晒的那个灰条菜和豇豆、扁豆、茄子干子、葫芦条儿，各样干 菜带些来――我们这里上上下下都爱吃这个――就算了。}
}
{
    \eA{Goody Liu gave utterance to profuse expressions of gratitude and signified
        her readiness to comply with her wishes. "Just you go to sleep," P'ing Erh
        urged, "and I'll get the things ready for you and put them in here. As soon
        as the day breaks to-morrow, I'll send the servant-lads to hire a cart and
        pack them in; don't you therefore worry yourself in the least on that
        score!"}
}
{
}

\Verse{59}
{
    \zA{别的一概不要，别罔费了 心。”}
}
{
    \eA{Goody Liu felt more and more ineffably grateful. So crossing over, she again
        said,}
}
{
}

\Verse{60}
{
    \zA{刘老老千恩万谢的答应了。}
}
{
    \eA{with warm protestations of thankfulness, good bye to lady Feng; after which,}
}
{
}

\Verse{61}
{
    \zA{平儿道：“你只管睡你的去，我替你收拾妥当了， 就放在这里，明儿一早打发小厮们雇辆车装上，不用你费一点心儿。”}
}
{
    \eA{she repaired to dowager lady Chia's quarters on this side, where she slept,
        with one sleep, during the whole night. Early the next day, as soon as she
        had combed her hair and performed her ablutions, she asked to go and pay her
        adieus to lady Chia.}
}
{
}

\Verse{62}
{
    \zA{刘老老越发 感激不尽，过来又千恩万谢的辞了凤姐儿，过贾母这边睡了一夜。}
}
{
    \eA{But as old lady Chia was unwell, the various members of the family came to
        see how she was getting on. On their reappearance outside, they transmitted
        orders that the doctor should be sent for.}
}
{
}

\Verse{63}
{
    \zA{次早梳洗了，就 要告辞。}
}
{
    \eA{In a little time, a matron reported that the doctor had arrived,}
}
{
}

\Verse{64}
{
    \zA{因贾母欠安，众人都过来请安，出去传请大夫。}
}
{
    \eA{and an old nurse invited dowager lady Chia to ensconce herself under the
        curtain. "I'm an old woman!" lady Chia remonstrated.}
}
{
}

\Verse{65}
{
    \zA{一时婆子回：“大夫来了。”}
}
{
    \eA{"Am I not aged enough to be a mother to that fellow? and am I, pray, to
        still stand on any ceremonies with him?}
}
{
}

\Verse{66}
{
    \zA{老 嬷嬷请贾母进幔子去坐，贾母道：“我也老了，那里养不出那阿物儿来，还怕他不 成!不用放幔子，就这样瞧罢。”}
}
{
    \eA{There's no need to drop the curtain; I'll see him as I am, and have done."
        Hearing her objections, the matrons fetched a small table, and, laying a
        small pillow on it, they directed a servant to ask the doctor in.}
}
{
}

\Verse{67}
{
    \zA{众婆子听了，便拿过一张小桌子来，放下一个小枕 头，便命人请。}
}
{
    \eA{Presently, they perceived the trio Chia Chen, Chia Lien, and Chia Jung,
        bringing Dr. Wang. Dr. Wang did not presume to use the raised road, but
        confining himself to the side steps,}
}
{
}

\Verse{68}
{
    \zA{一时只见贾珍、贾琏、贾蓉三个人，将王太医领来。}
}
{
    \eA{he kept pace with Chia Chen until they reached the platform. Two matrons,
        who had been standing, one on either side from an early hour,}
}
{
}

\Verse{69}
{
    \zA{王太医不敢走 甬路，只走旁阶，跟着贾珍到了台阶上。}
}
{
    \eA{raised the portiére. A couple of old women servants then took the lead and
        showed the way in. But Pao-yü too appeared on the scene to meet them.}
}
{
}

\Verse{70}
{
    \zA{早有两个婆子在两边打起帘子，两个婆子 在前导引进去，又见宝玉迎接出来。}
}
{
    \eA{They found old lady Chia seated bolt upright on the couch, dressed in a blue
        crape jacket, lined with sheep skin, every curl of which resembled a pearl.
        On the right and left stood four young maids, whose hair had not as yet been
        allowed to grow,}
}
{
}

\Verse{71}
{
    \zA{见贾母穿着青绉绸一斗珠儿的羊皮褂子，端坐 在榻上。}
}
{
    \eA{with fly-brushes, finger-bowls, and other such articles in their hands. Five
        or six old nurses were also drawn up on both sides like wings.}
}
{
}

\Verse{72}
{
    \zA{两边四个未留头的小丫鬟，都拿着蝇刷漱盂等物，又有五六个老嬷嬷雁翅 摆在两旁。}
}
{
    \eA{At the back of the jade-green gauze mosquito-house were faintly visible
        several persons in red and green habiliments, with gems on their heads, and
        gold trinkets in their coiffures. Dr. Wang could not muster the courage to
        raise his head.}
}
{
}

\Verse{73}
{
    \zA{碧纱厨后，隐隐约约有许多穿红着绿、戴宝插金的人，王太医也不敢抬 头，忙上来请了安。}
}
{
    \eA{With speedy step, he advanced and paid his obeisance. Dowager lady Chia
        noticed that he wore the official dress of the sixth grade, and she
        accordingly concluded that he must be an imperial physician. "How are you
        noble doctor?"}
}
{
}

\Verse{74}
{
    \zA{贾母见他穿着六品服色，便知是御医了，含笑问：“供奉好？”}
}
{
    \eA{she inquired, forcing a smile. "What is the worthy surname of this noble
        doctor?" she then asked Chia Chen. Chia Chen and his companions made prompt
        reply.}
}
{
}

\Verse{75}
{
    \zA{因问贾珍：“这位供奉贵姓？”}
}
{
    \eA{"His surname is Wang," they said. "There was once a certain Wang Chün-hsiao
        who filled the chair of President of the College of Imperial Physicians,"}
}
{
}

\Verse{76}
{
    \zA{贾珍等忙回：“姓王。”}
}
{
    \eA{dowager lady smilingly proceeded. "He excelled in feeling the pulse."}
}
{
}

\Verse{77}
{
    \zA{贾母笑道：“当日太医院正堂 有个王君效，好脉息。”}
}
{
    \eA{Dr. Wang bent his body, and with alacrity he lowered his head and returned
        her smile. "That was," he explained, "my grand uncle." "Is it really so!"}
}
{
}

\Verse{78}
{
    \zA{王太医忙躬身低头含笑，因说：“那是晚生家叔祖。”}
}
{
    \eA{laughingly pursued dowager lady Chia, upon catching this reply. "We can then
        call ourselves old friends!" So speaking, she quietly put out her hand and
        rested it on the small pillow.}
}
{
}

\Verse{79}
{
    \zA{贾母 听了笑道：“原来这样，也算是世交了。”}
}
{
    \eA{A nurse laid hold of a small stool and placed it before the small table,
        slightly to the side of it. Dr. Wang bent one knee and took a seat on the
        stool.}
}
{
}

\Verse{80}
{
    \zA{一面说，一面慢慢的伸手放在小枕头上。}
}
{
    \eA{Drooping his head, he felt the pulse of the one hand for a long while; next,
        he examined that of the other;}
}
{
}

\Verse{81}
{
    \zA{嬷嬷端着一张小杌子放在小桌前面，略偏些。}
}
{
    \eA{after which, hastily making a curtsey, he bent his head and started on his
        way out of the apartment. "Excuse me for the trouble I've put you to!"}
}
{
}

\Verse{82}
{
    \zA{王太医便盘着一条腿儿坐下，歪着头 诊了半日，又诊了那只手，忙欠身低头退出。}
}
{
    \eA{dowager lady Chia smiled. "Chen Erh, escort him outside, and do see that he
        has a cup of tea." Chia Chen, Chia Lien and the rest of their companions
        immediately acquiesced by uttering several yes's,}
}
{
}

\Verse{83}
{
    \zA{贾母笑说：“劳动了。}
}
{
    \eA{and once more they led Dr. Wang into the outer study.}
}
{
}

\Verse{84}
{
    \zA{珍哥让出去， 好生看茶。”}
}
{
    \eA{"Your worthy senior," Dr. Wang explained, "has nothing else the matter with
        her than a slight chill,}
}
{
}

\Verse{85}
{
    \zA{贾珍、贾琏等忙答应了几个“是”，复领王太医到外书房中。}
}
{
    \eA{which she must have inadvertently contracted. She needn't, after all, take
        any medicines; all she need do is to diet herself and keep warm a little;
        and she'll get all right.}
}
{
}

\Verse{86}
{
    \zA{王太医说： “太夫人并无别症，偶感了些风寒，其实不用吃药，不过略清淡些，常暖着点儿， 就好了。}
}
{
    \eA{But I'll now write a prescription, in here. Should her venerable ladyship
        care to take any of the medicine, then prepare a dose, according to the
        prescription, and let her have it. But should she be loth to have any, well,
        never mind, it won't be of any consequence."}
}
{
}

\Verse{87}
{
    \zA{如今写个方子在这里，若老人家爱吃，便按方煎一剂吃；若懒怠吃，也就 罢了。”}
}
{
    \eA{Saying this, he wrote the prescription, as he sipped his tea. But when about
        to take his leave, he saw a nurse bring Ta Chieh-erh into the room. "Mr.
        Wang," she said, "do also have a look at our Chieh Erh!"}
}
{
}

\Verse{88}
{
    \zA{说着，吃茶，写了方子。}
}
{
    \eA{Upon hearing her appeal, Dr. Wang immediately rose to his feet.}
}
{
}

\Verse{89}
{
    \zA{刚要告辞，只见奶子抱了大姐儿出来，笑说：“王 老爷也瞧瞧我们。”}
}
{
    \eA{While she was clasped in her nurse's arms, he rested Ta Chieh-erh's hand on
        his left hand and felt her pulse with his right, and rubbing her forehead,
        he asked her to put out her tongue and let him see it.}
}
{
}

\Verse{90}
{
    \zA{王太医听说，忙起身就奶子怀中，左手托着大姐儿的手，右手 诊了一诊，又摸了一摸头，又叫伸出舌头来瞧瞧，笑道：“我要说了，妞儿该骂我
        了：只要清清净净的饿两顿就好了。}
}
{
    \eA{"Were I to express my views about Chieh Erh, you would again abuse me! If
        she's, however, kept quiet and allowed to go hungry for a couple of meals,
        she'll get over this. There's no necessity for her to take any decocted
        medicines. I'll just send her some pills,}
}
{
}

\Verse{91}
{
    \zA{不必吃煎药，我送点丸药来，临睡用姜汤研开 吃下去就好了。”}
}
{
    \eA{which you'll have to dissolve in a preparation of ginger, and give them to
        her before she goes to sleep; when she has had these, there will be nothing
        more the matter with her."}
}
{
}

\Verse{92}
{
    \zA{说毕，告辞而去。}
}
{
    \eA{At the conclusion of these recommendations,}
}
{
}

\Verse{93}
{
    \zA{贾珍等拿了药方来回明贾母原故，将药方放在 案上出去，不在话下。}
}
{
    \eA{he bade them goodbye and took his departure. Chia Chen and his companions
        then took the prescription and came and explained to old lady Chia the
        nature of her indisposition, and, depositing on the table,}
}
{
}

\Verse{94}
{
    \zA{这里王夫人和李纨、凤姐儿、宝钗姐妹等，见大夫出去，方从厨后出来。}
}
{
    \eA{the paper given to them by the doctor, they quitted her presence. But
        nothing more need be said about them. Madame Wang and Li Wan, lady Feng, Pao
        Ch'ai and the other young ladies noticed,}
}
{
}

\Verse{95}
{
    \zA{王夫 人略坐一坐，也回房去了。}
}
{
    \eA{meanwhile, that the doctor had gone, and they eventually egressed from the
        back of the mosquito-house.}
}
{
}

\Verse{96}
{
    \zA{刘老老见无事，方上来和贾母告辞。}
}
{
    \eA{After a short stay, Madame Wang returned to her quarters. Goody Liu
        repaired, when she perceived everything quiet again,}
}
{
}

\Verse{97}
{
    \zA{贾母说：“闲了再 来。”}
}
{
    \eA{into the upper rooms and made her adieus to dowager lady Chia.}
}
{
}

\Verse{98}
{
    \zA{又命鸳鸯来：“好生打发刘老老出去。}
}
{
    \eA{"When you've got any leisure, do pay us another visit," old lady Chia urged,
        and bidding Yuan Yang come to her,}
}
{
}

\Verse{99}
{
    \zA{我身上不好，不能送你。”}
}
{
    \eA{"Do be careful," she added, "and see dame Liu safely on her way out;}
}
{
}

\Verse{100}
{
    \zA{刘老老道了 谢，又作辞，方同鸳鸯出来。}
}
{
    \eA{for not being well I can't escort you myself." Goody Liu expressed her
        thanks, and saying good bye a second time,}
}
{
}

\Verse{101}
{
    \zA{到了下房，鸳鸯指炕上一个包袱说道：“这是老太太 的几件衣裳，都是往年间生日节下众人孝敬的。}
}
{
    \eA{she betook herself, along with Yüan Yang, into the servants' quarters. Here
        Yüan Yang pointed at a bundle on the stove-couch. "These are," she said,
        "several articles of clothing, belonging to our old mistress; they were
        presented to her in years gone by,}
}
{
}

\Verse{102}
{
    \zA{老太太从不穿人家做的，收着也可 惜，却是一次也没穿过的，昨日叫我拿出两套来送你带了去，或送人，或自己家里 穿罢。}
}
{
    \eA{by members of our family on her birthdays and various festivals; her
        ladyship never wears anything made by people outside; yet to hoard these
        would be a downright pity! Indeed, she hasn't worn them even once. It was
        yesterday that she told me to get out two costumes and hand them to you to
        take along with you,}
}
{
}

\Verse{103}
{
    \zA{这盒子里头是你要的面果子。}
}
{
    \eA{either to give as presents, or to be worn by some one in your home; but
        don't make fun of us!}
}
{
}

\Verse{104}
{
    \zA{这包儿里头是你前儿说的药，梅花点舌丹也有， 紫金锭也有，活络丹也有，催生保命丹也有：每一样是一张方子包着，总包在里头 了。}
}
{
    \eA{In the box you'll find the flour-fruits, for which you asked. This bundle
        contains the medicines to which you alluded the other day. There are
        'plum-blossom-spotted-tongue pills,' and 'purple-gold- ingot- pills,' also
        'vivifying-blood-vessels-pills,'}
}
{
}

\Verse{105}
{
    \zA{这是两个荷包，带着玩罢。”}
}
{
    \eA{as well as 'driving-offspring and preserving-life pills;'}
}
{
}

\Verse{106}
{
    \zA{说着，又抽开系子，掏出两个“笔锭如意”的锞 子来给他瞧，又笑道：“荷包你拿去，这个留下给我罢。”}
}
{
    \eA{each kind being rolled up in a sheet bearing the prescription; and the whole
        lot of them are packed up in here. While these two are purses for you to
        wear in the way of ornaments." So saying, she forthwith loosened the cord,
        and, producing two ingots representing pencils,}
}
{
}

\Verse{107}
{
    \zA{刘老老已喜出望外，早又 念了几千佛，听鸳鸯如此说，便忙说道：“姑娘只管留下罢。”}
}
{
    \eA{and with 'ju i' on them, implying 'your wishes will surely be fulfilled,'
        she drew near and showed them to her, "Take the purses," she pursued
        smiling, "but do leave these behind and give them to me."}
}
{
}

\Verse{108}
{
    \zA{鸳鸯见他信以为真， 笑着仍给他装上，说道：“哄你玩呢!我有好些呢。}
}
{
    \eA{Goody Liu was so overjoyed that she had, from an early period, come out
        afresh with several thousands of invocations of Buddha's names. When she
        therefore heard Yüan Yang's suggestion,}
}
{
}

\Verse{109}
{
    \zA{留着年下给小孩子们罢。”}
}
{
    \eA{"Miss," she quickly rejoined, "you're at perfect liberty to keep them!"}
}
{
}

\Verse{110}
{
    \zA{说着， 只见一个小丫头拿着个成窑钟子来，递给刘老老，说：“这是宝二爷给你的。”}
}
{
    \eA{Yüan Yang perceived that her words were believed by her; so smiling she once
        more dropped the ingots into the purse. "I was only joking with you for
        fun!" she observed. "I've got a good many like these;}
}
{
}

\Verse{111}
{
    \zA{刘老 老道：“这是那里说起？}
}
{
    \eA{keep them therefore and give them, at the close of the year, to your young
        children."}
}
{
}

\Verse{112}
{
    \zA{我那一世修来的，今儿这样！”}
}
{
    \eA{Speaking the while, she espied a young maid walk in with a cup from the
        'Ch'eng' kiln, and hand it to old goody Liu.}
}
{
}

\Verse{113}
{
    \zA{说着便接过来。}
}
{
    \eA{"This," (she said,) "our master Secundus,}
}
{
}

\Verse{114}
{
    \zA{鸳鸯道：“前 儿我叫你洗澡，换的衣裳是我的，你不弃嫌，我还有几件也送你罢。”}
}
{
    \eA{Mr. Pao, gives you." "Whence could I begin enumerating the things I got!"
        Goody Liu exclaimed. "In what previous existence did I accomplish anything
        so meritorious as to bring to-day this heap of blessings upon me!" With
        these words, she eagerly took possession of the cup. "The clothes I gave you
        the other day,}
}
{
}

\Verse{115}
{
    \zA{刘老老又忙 道谢。}
}
{
    \eA{when I asked you to have a bath, were my own," Yüan Yang resumed,}
}
{
}

\Verse{116}
{
    \zA{鸳鸯果然又拿出几件来，给他包好。}
}
{
    \eA{"and if you don't think them too mean, I've got a few more, which I would
        also like to let you have." Goody Liu thanked her with vehemence,}
}
{
}

\Verse{117}
{
    \zA{刘老老又要到园中辞谢宝玉和众姊妹王 夫人等去，鸳鸯道：“不用去了。}
}
{
    \eA{so Yüan Yang, in point of fact, produced several more articles of clothing,
        and these she packed up for her. Goody Liu thereupon expressed a desire to
        also go into the garden and take leave of Pao-yü and the young ladies,}
}
{
}

\Verse{118}
{
    \zA{他们这会子也不见人，回来我替你说罢。}
}
{
    \eA{Madame Wang and the other inmates and to thank them for all they did for
        her, but Yüan Yang raised objections.}
}
{
}

\Verse{119}
{
    \zA{闲了再 来。”}
}
{
    \eA{"You can dispense with going!"}
}
{
}

\Verse{120}
{
    \zA{又命了一个老婆子，吩咐他：“二门上叫两个小厮来，帮着老老拿了东西送去。”}
}
{
    \eA{she remarked. "They don't see any one just now! But I'll deliver the message
        for you by and bye! When you've got any leisure, do come again. Go to the
        second gate," she went on to direct an old matron,}
}
{
}

\Verse{121}
{
    \zA{婆子答应了。}
}
{
    \eA{"and call two servant-lads to come here,}
}
{
}

\Verse{122}
{
    \zA{又和刘老老到了凤姐儿那边，一并拿了东西，在角门上命小厮们搬出 去，直送刘老老上车去了，不在话下。}
}
{
    \eA{and help this old dame to take her things away!" After the matron had
        signified her obedience, Yüan Yang returned with goody Liu to lady Feng's
        quarters, on the off part of the mansion, and, taking the presents as far as
        the side gate, she bade the servant-lads carry them out.}
}
{
}

\Verse{123}
{
    \zA{且说宝钗等吃过早饭，又往贾母处问安，回园至分路之处，宝钗便叫黛玉道： “颦儿跟我来!有一句话问你。”}
}
{
    \eA{She herself then saw goody Liu into her curricle and start on her journey
        homewards. But without commenting further on this topic, let us revert to
        Pao-ch'ai and the other girls. After breakfast, they recrossed into their
        grandmother's rooms and made inquiries about her health. On their way back
        to the garden,}
}
{
}

\Verse{124}
{
    \zA{黛玉便笑着跟了来。}
}
{
    \eA{they reached a point where they had to take different roads.}
}
{
}

\Verse{125}
{
    \zA{至蘅芜院中，进了房，宝钗便 坐下，笑道：“你还不给我跪下!我要审你呢。”}
}
{
    \eA{Pao-ch'ai then called out to Tai-yü. "P'in Erh!" she observed, "come with
        me; I've got a question to ask you." Tai-yü wended her steps therefore with
        Pao-ch'ai into the Heng Wu court.}
}
{
}

\Verse{126}
{
    \zA{黛玉不解何故，因笑道：“你瞧宝丫 头疯了!审我什么？”}
}
{
    \eA{As soon as they entered the house, Pao-ch'ai threw herself into a seat.
        "Kneel down!" she smiled. "I want to examine you about something!"}
}
{
}

\Verse{127}
{
    \zA{宝钗冷笑道：“好个千金小姐!好个不出屋门的女孩儿!满嘴里 说的是什么？}
}
{
    \eA{Tai-yü could not fathom her object, and consequently laughed. "Look here."
        she cried, "this chit Pao has gone clean off her senses! What do you want to
        examine me about?" Pao-ch'ai gave a sardonic smile.}
}
{
}

\Verse{128}
{
    \zA{你只实说罢。”}
}
{
    \eA{"My dear, precious girl, my dear maiden," she exclaimed,}
}
{
}

\Verse{129}
{
    \zA{黛玉不解，只管发笑，心里也不免疑惑，口里只说：“我 何曾说什么？}
}
{
    \eA{"what utter trash fills your mouth! Just speak the honest and candid truth,
        and finish!" Tai-yü could so little guess her meaning that her sole resource
        was to smile.}
}
{
}

\Verse{130}
{
    \zA{你不过要捏我的错儿罢咧。}
}
{
    \eA{Inwardly, however, she could not help beginning to experience certain
        misgivings.}
}
{
}

\Verse{131}
{
    \zA{你倒说出来我听听。”}
}
{
    \eA{"What did I say?" she remarked. "You're bent upon picking out my faults!}
}
{
}

\Verse{132}
{
    \zA{宝钗笑道：“你还装 憨儿呢!昨儿行酒令儿，你说的是什么？}
}
{
    \eA{Speak out and let me hear what it's all about!" "Do you still pretend to be
        a fool?" Pao-ch'ai laughed. "When we played yesterday that game of
        wine-forfeits,}
}
{
}

\Verse{133}
{
    \zA{我竟不知是那里来的。”}
}
{
    \eA{what did you say? I really couldn't make out any head or tail."}
}
{
}

\Verse{134}
{
    \zA{黛玉一想，方想起 昨儿失于检点，那《牡丹亭》、《西厢记》说了两句，不觉红了脸，便上来搂着宝钗 笑道：“好姐姐!原是我不知道，随口说的。}
}
{
    \eA{Tai-yü, after a moment's reflection, remembered eventually that she had the
        previous day been guilty of a slip of the tongue and come out with a couple
        of passages from the 'Peony Pavilion,' and the 'Record of the West
        Side-house,' and, of a sudden, her face got scarlet with blushes. Drawing
        near Pao-ch'ai she threw her arms round her. "My dear cousin!" she smiled,}
}
{
}

\Verse{135}
{
    \zA{你教给我，再不说了。”}
}
{
    \eA{"I really wasn't conscious of what I was saying! It just blurted out of my
        mouth!}
}
{
}

\Verse{136}
{
    \zA{宝钗笑道：“我 也不知道，听你说的怪好的，所以请教你。”}
}
{
    \eA{But now that you've called me to task, I won't say such things again." "I've
        no idea of what you were driving at," Pao-ch'ai laughingly rejoined.}
}
{
}

\Verse{137}
{
    \zA{黛玉道：“好姐姐!你别说给别人，我 再不说了！”}
}
{
    \eA{"What I heard you recite sounds so thoroughly unfamiliar to me, that I beg
        you to enlighten me!" "Dear cousin," pleaded Tai-yü,}
}
{
}

\Verse{138}
{
    \zA{宝钗见他羞的满脸飞红，满口央告，便不肯再往下问。}
}
{
    \eA{"don't tell anyone else! I won't, in the future, breathe such things again."
        Pao-ch'ai noticed how from shame the blood rushed to her face,}
}
{
}

\Verse{139}
{
    \zA{因拉他坐下吃 茶，款款的告诉他道：“你当我是谁？}
}
{
    \eA{and how vehement she was in her entreaties, and she felt both to press her
        with questions; so pulling her into a seat to make her have a cup of tea,}
}
{
}

\Verse{140}
{
    \zA{我也是个淘气的，从小儿七八岁上，也够个人 缠的。}
}
{
    \eA{she said to her in a gentle tone, "Whom do you take me for? I too am
        wayward; from my youth up, yea ever since I was seven or eight,}
}
{
}

\Verse{141}
{
    \zA{我们家也算是个读书人家，祖父手里也极爱藏书。}
}
{
    \eA{I've been enough trouble to people! Our family was also what one would term
        literary. My grandfather's extreme delight was to be ever with a book in his
        hand.}
}
{
}

\Verse{142}
{
    \zA{先时人口多，姐妹弟兄也 在一处，都怕看正经书。}
}
{
    \eA{At one time, we numbered many members, and sisters and brothers all lived
        together; but we had a distaste for wholesome books.}
}
{
}

\Verse{143}
{
    \zA{弟兄们也有爱诗的，也有爱词的，诸如这些《西厢》、《琵 琶》以及《元人百种》，无所不有。}
}
{
    \eA{Among my brothers, some were partial to verses; others had a weakness for
        blank poetical compositions; and there were none of such works as the
        'Western side-House,' and 'the Guitar,' even up to the hundred and one books
        of the 'Yüan' authors,}
}
{
}

\Verse{144}
{
    \zA{他们背着我们偷看，我们也背着他们偷看。}
}
{
    \eA{which they hadn't managed to get. These books they stealthily read behind
        our backs; but we, on our part, devoured them,}
}
{
}

\Verse{145}
{
    \zA{后 来大人知道了，打的打，骂的骂，烧的烧，丢开了。}
}
{
    \eA{on the sly, without their knowing it. Subsequently, our father came to get
        wind of it; and some of us he beat, while others he scolded;}
}
{
}

\Verse{146}
{
    \zA{所以咱们女孩儿家不认字的倒 好：男人们读书不明理，尚且不如不读书的好，何况你我？}
}
{
    \eA{burning some of the books, and throwing away others. It is therefore as well
        that we girls shouldn't know anything of letters. Men, who study books and
        don't understand the right principle, can't, moreover, reach the standard of
        those,}
}
{
}

\Verse{147}
{
    \zA{连做诗写字等事，这也 不是你我分内之事，究竟也不是男人分内之事。}
}
{
    \eA{who don't go in for books; so how much more such as ourselves? Even
        versifying, writing and the like pursuits aren't in the line of such as you
        and me. Indeed, neither are they within the portion of men.}
}
{
}

\Verse{148}
{
    \zA{男人们读书明理，辅国治民，这才 是好。}
}
{
    \eA{Men, who go in for study and fathom the right principles, should cooperate
        in the government of the empire,}
}
{
}

\Verse{149}
{
    \zA{只是如今并听不见有这样的人，读了书，倒更坏了。}
}
{
    \eA{and should rule the nation; this would be a nobler purpose; but one doesn't
        now-a-days hear of the very existence of such persons!}
}
{
}

\Verse{150}
{
    \zA{这并不是书误了他，可 惜他把书遭塌了，所以竟不如耕种买卖，倒没有什么大害处。}
}
{
    \eA{Hence, the study of books makes them worse than they ever were before. But
        it isn't the books that ruin them; the misfortune is that they make improper
        use of books! That is why study doesn't come up to ploughing and sowing and
        trading;}
}
{
}

\Verse{151}
{
    \zA{至于你我，只该做些 针线纺绩的事才是；偏又认得几个字。}
}
{
    \eA{as these pursuits exercise no serious pernicious influences. As far,
        however, as you and I go, we should devote our minds simply to matters
        connected with needlework and spinning;}
}
{
}

\Verse{152}
{
    \zA{既认得了字，不过拣那正经书看也罢了，最 怕见些杂书，移了性情，就不可救了。”}
}
{
    \eA{for we will then be fulfilling our legitimate duties. Yet, it so happens
        that we too know a few characters. But, as we can read, it behoves us to
        choose no other than wholesome works; for these will do us no harm! What are
        most to be shirked are those low books,}
}
{
}

\Verse{153}
{
    \zA{一席话，说的黛玉垂头吃茶，心下暗服， 只有答应“是”的一字。}
}
{
    \eA{as, when once they pervert the disposition, there remains no remedy
        whatever!" While she indulged in this long rigmarole, Tai-yü lowered her
        head and sipped her tea.}
}
{
}

\Verse{154}
{
    \zA{忽见素云进来说：“我们奶奶请二位姑娘商议要紧的事呢。}
}
{
    \eA{And though she secretly shared the same views on the subject, all the answer
        she gave her in assent was limited to one single word 'yes.' But at an
        unexpected moment,}
}
{
}

\Verse{155}
{
    \zA{二姑娘、三姑娘、 四姑娘、史姑娘、宝二爷，都等着呢。”}
}
{
    \eA{Su Yün appeared in the room. "Our lady Lien," she said, "requests the
        presence of both of you, young ladies, to consult with you in an important
        matter.}
}
{
}

\Verse{156}
{
    \zA{宝钗说：“又是什么事？”}
}
{
    \eA{Miss Secunda, Miss Tertia, Miss Quarta, Miss Shih and Mr. Pao,}
}
{
}

\Verse{157}
{
    \zA{黛玉道：“咱们 到了那里就知道了。”}
}
{
    \eA{our master Secundus, are there waiting for you." "What's up again?"
        Pao-ch'ai inquired.}
}
{
}

\Verse{158}
{
    \zA{说着，便和宝钗往稻香村来，果见众人都在那里。}
}
{
    \eA{"You and I will know what it is when we get there," Tai-yü explained. So
        saying, she came, with Pao-ch'ai, into the Tao Hsiang village.}
}
{
}

\Verse{159}
{
    \zA{李纨见了 他两个，笑道：“社还没起，就有脱滑儿的了，四丫头要告一年的假呢。”}
}
{
    \eA{Here they, in fact, discovered every one assembled. As soon as Li Wan caught
        sight of the two cousins, she smiled. "The society has barely been started,"
        she observed, "and here's one who wants to give us the slip;}
}
{
}

\Verse{160}
{
    \zA{黛玉笑道： “都是老太太昨儿一句话，又叫他画什么园子图儿，惹的他乐得告假了。”}
}
{
    \eA{that girl Quarta wishes to apply for a whole year's leave." "It's that
        single remark of our worthy senior's yesterday that is at the bottom of it!"
        Tai-yü laughed. "For by bidding her execute some painting or other of the
        garden,}
}
{
}

\Verse{161}
{
    \zA{探春笑 道：“也别怪老太太，都是刘老老一句话。”}
}
{
    \eA{she has put her in such high feather that she applies for leave!" "Don't be
        so hard upon our dear ancestor!" Pao-Ch'ai rejoined, a smile playing on her
        lips.}
}
{
}

\Verse{162}
{
    \zA{黛玉忙笑接道：“可是呢，都是他一句 话。}
}
{
    \eA{"It's entirely due to that allusion of grandmother Liu's." Tai-yü speedily
        took up the thread of the conversation.}
}
{
}

\Verse{163}
{
    \zA{他是那一门子的老老？}
}
{
    \eA{"Quite so!" she smiled. "It's all through that remark of hers!}
}
{
}

\Verse{164}
{
    \zA{直叫他是个‘母蝗虫’就是了。”}
}
{
    \eA{But of what branch of the family is she a grandmother? We should merely
        address her as the 'female locust;'}
}
{
}

\Verse{165}
{
    \zA{说着，大家都笑起来。}
}
{
    \eA{that's all." As she spoke, one and all were highly amused.}
}
{
}

\Verse{166}
{
    \zA{宝钗笑道：“世上的话，到了二嫂子嘴里也就尽了，幸而二嫂子不认得字，不大通， 不过一概是市俗取笑儿。}
}
{
    \eA{"When any mortal language finds its way into that girl Feng's mouth,"
        Pao-ch'ai laughed, "she knows how to turn it to the best account! What a
        fortunate thing it is that that vixen Feng has no idea of letters and can't
        boast of much culture! Her \_forte\_ is simply such vulgar things as suffice
        to raise a laugh!}
}
{
}

\Verse{167}
{
    \zA{更有颦儿这促狭嘴，他用《春秋》的法子，把市俗粗话撮 其要，删其繁，再加润色，比方出来，一句是一句。}
}
{
    \eA{Worse than her is that P'in Erh with that coarse tongue! She has recourse to
        the devices of the 'Ch'un Ch'iu'! By selecting, from the vulgar expressions
        used in low slang, the most noteworthy points, she eliminates what's
        commonplace, and makes, with the addition of a little elegance and finish,}
}
{
}

\Verse{168}
{
    \zA{这‘母蝗虫’三字，把昨儿那 些形景都画出来了。}
}
{
    \eA{her style so much like that of the text that each sentence has a peculiar
        character of its own! The three words representing 'female locust' bring out
        clearly the various circumstances connected with yesterday!}
}
{
}

\Verse{169}
{
    \zA{亏他想的倒也快！”}
}
{
    \eA{The wonder is that she has been so quick in devising them!"}
}
{
}

\Verse{170}
{
    \zA{众人听了，都笑道：“你这一注解，也就不 在他两个以下了。”}
}
{
    \eA{After lending an ear to her arguments, they all laughed. "Those explanations
        of yours," they cried, "show well enough that you are not below those two!"}
}
{
}

\Verse{171}
{
    \zA{李纨道：“我请你们大家商议，给他多少日子的假？}
}
{
    \eA{"Pray, let's consult as to how many days' leave to grant her!" Li Wan
        proposed. "I gave her a month, but she thinks it too little.}
}
{
}

\Verse{172}
{
    \zA{我给了他一个月的假，他嫌 少，你们怎么说？”}
}
{
    \eA{What do you say about it?" "Properly speaking," Tai-yü put in, "one year
        isn't much! The laying out of this garden occupied a whole year;}
}
{
}

\Verse{173}
{
    \zA{黛玉道：“论理，一年也不多，这园子盖就盖了一年，如今要 画，自然得二年的工夫呢：又要研墨，又要蘸笔，又要铺纸，又要着颜色，又要―
        ―”刚说到这里，黛玉也自己掌不住，笑道：“又要照着样儿慢慢的画，可不得二 年的工夫？”}
}
{
    \eA{and to paint a picture of it now will certainly need two years' time. She'll
        have to rub the ink, to moisten the pencils, to stretch the paper, to mix
        the pigments, and to…." When she had reached this point, even Tai-yü could
        not restrain herself from laughing.}
}
{
}

\Verse{174}
{
    \zA{众人听了，都拍手笑个不住。}
}
{
    \eA{"If she goes on so leisurely to work," she exclaimed, "won't she require two
        years' time?"}
}
{
}

\Verse{175}
{
    \zA{宝钗笑道：“有趣!最妙落后一句是‘慢 慢的画’。}
}
{
    \eA{Those, who caught this insinuation, clapped their hands and indulged in
        incessant merriment. "Her innuendoes are full of zest!"}
}
{
}

\Verse{176}
{
    \zA{他可不画去，怎么就有了呢？}
}
{
    \eA{Pao-ch'ai ventured laughingly. "But what takes the cake is that last remark
        about leisurely going to work,}
}
{
}

\Verse{177}
{
    \zA{所以昨儿那些笑话儿虽然可笑，回想是没趣 的。}
}
{
    \eA{for if she weren't to paint at all, how could she ever finish her task?
        Hence those jokes cracked yesterday were,}
}
{
}

\Verse{178}
{
    \zA{你们细想，颦儿这几句话，虽没什么，回想却有滋味。}
}
{
    \eA{sufficient, of course, to evoke laughter, but, on second thought, they're
        devoid of any fun! Just you carefully ponder over P'in Erh's words!}
}
{
}

\Verse{179}
{
    \zA{我倒笑的动不得了。”}
}
{
    \eA{Albeit they don't amount to much, you'll nevertheless find,}
}
{
}

\Verse{180}
{
    \zA{惜春道：“都是宝姐姐赞的他越发逞强，这会子又拿我取笑儿。”}
}
{
    \eA{when you come to reflect on them, that there's plenty of gusto about them.
        I've really had such a laugh over them that I can scarcely move! "It's the
        way that cousin Pao-ch'ai puffs her up,"}
}
{
}

\Verse{181}
{
    \zA{黛玉忙拉他笑道： “我且问你，还是单画这园子呢，还是连我们众人都画在上头呢？”}
}
{
    \eA{Hsi Ch'un observed "that makes her so much the more arrogant that she turns
        me also into a laughing-stock now!" Tai-yü hastily smiled and pulled her
        towards her. "Let me ask you," she said, "are you only going to paint the
        garden,}
}
{
}

\Verse{182}
{
    \zA{惜春道：“原 是只画这园子。}
}
{
    \eA{or will you insert us in it as well?" "My original idea was to have simply
        painted the garden,"}
}
{
}

\Verse{183}
{
    \zA{昨儿老太太又说：‘单画园子，成了房样子了。’}
}
{
    \eA{Hsi Ch'un explained; "but our worthy senior told me again yesterday that a
        mere picture of the grounds would resemble the plan of a house, and
        recommended that I should introduce some inmates too so as to make it look
        like what a painting should.}
}
{
}

\Verse{184}
{
    \zA{叫连人都画上，就 像行乐图儿才好。}
}
{
    \eA{I've neither the knack for the fine work necessary for towers and terraces,
        nor have I the skill to draw representations of human beings;}
}
{
}

\Verse{185}
{
    \zA{我又不会这工细楼台，又不会画人物，又不好驳回，正为这个为 难呢。”}
}
{
    \eA{but as I couldn't very well raise any objections, I find myself at present
        on the horns of a dilemma about it!" "Human beings are an easy matter!"
        Tai-yü said. "What beats you are insects."}
}
{
}

\Verse{186}
{
    \zA{黛玉道：“人物还容易，你草虫儿上不能。”}
}
{
    \eA{"Here you are again with your trash!" Li Wan exclaimed. "Will there be any
        need to also introduce insects in it?}
}
{
}

\Verse{187}
{
    \zA{李纨道：“你又说不通的话了。}
}
{
    \eA{As far, however, as birds go, it may probably be advisable to introduce one
        or two kinds!"}
}
{
}

\Verse{188}
{
    \zA{这上头那里又用草虫儿呢？}
}
{
    \eA{"If any other insects are not put in the picture," Tai-yü smiled, "it won't
        matter;}
}
{
}

\Verse{189}
{
    \zA{或者翎毛倒要点缀一两样。”}
}
{
    \eA{but without yesterday's female locust in it, it will fall short of the
        original?"}
}
{
}

\Verse{190}
{
    \zA{黛玉笑道：“别的草虫儿罢 了，昨儿的‘母蝗虫’不画上，岂不缺了典呢？”}
}
{
    \eA{This retort evoked further general amusement. While Tai-yü laughed, she beat
        her chest with both hands. "Begin painting at once!" she cried. "I've even
        got the title all ready.}
}
{
}

\Verse{191}
{
    \zA{众人听了，都笑起来。}
}
{
    \eA{The name I've chosen is, 'Picture of a locust brought in to have a good
        feed.'"}
}
{
}

\Verse{192}
{
    \zA{黛玉一面 笑的两只手捧着胸口，一面说道：“你快画罢，我连题跋都有了：起了名字，就叫 做《携蝗大嚼图》。”}
}
{
    \eA{At these words, they laughed so much the more heartily that at a time they
        bent forward, and at another they leant back. But a sound of "Ku tung" then
        fell on their ears, and unable to make out what could have dropped, they
        anxiously and precipitately looked about. It was, they found, Shih
        Hsiang-yün,}
}
{
}

\Verse{193}
{
    \zA{众人听了，越发哄然大笑的前仰后合。}
}
{
    \eA{who had been reclining on the back of the chair. The chair had, from the
        very outset, not been put in a sure place,}
}
{
}

\Verse{194}
{
    \zA{只听咕咚一声响，不知 什么倒了，急忙看时，原来是湘云伏在椅子背儿上，那椅子原不曾放稳，被他全身
        伏着背子大笑，他又不防，两下里错了笋，向东一歪，连人带椅子都歪倒了。}
}
{
    \eA{and while indulging in hearty merriment she threw her whole weight on the
        back. She did not, besides, notice that the dovetails on each side had come
        out, so with a tilt towards the east, she as well as the chair toppled over
        in a heap. Luckily, the wooden partition-wall was close enough to arrest her
        fall, and she did not sprawl on the ground. The sight of her created more
        amusement than ever among all her relatives;}
}
{
}

\Verse{195}
{
    \zA{幸有 板壁挡住，不曾落地。}
}
{
    \eA{so much so, that they could scarcely regain their equilibrium. It was only
        after Pao-yü had rushed up to her,}
}
{
}

\Verse{196}
{
    \zA{众人一见，越发笑个不住。}
}
{
    \eA{and given her a hand and raised her to her feet again that they at last
        managed to gradually stop laughing. Pao-yü then winked at Tai-yü.}
}
{
}

\Verse{197}
{
    \zA{宝玉忙赶上去扶住了起来，方渐 渐止了笑。}
}
{
    \eA{Tai-yü grasped his meaning, and, forthwith withdrawing into the inner room,
        she lifted the cover of the mirror,}
}
{
}

\Verse{198}
{
    \zA{宝玉和黛玉使个眼色儿，黛玉会意，便走至里间，将镜袱揭起。}
}
{
    \eA{and looked at her face. She found the hair about her temples slightly
        dishevelled, so, promptly opening Li Wan's toilet-case, and extracting a
        narrow brush, she stood in front of the mirror,}
}
{
}

\Verse{199}
{
    \zA{照了照，只见 两鬓略松了些，忙开了李纨的妆奁，拿出抿子来，对镜抿了两抿，仍旧收拾好了，
        方出来指着李纨道：“这是叫你带着我们做针线、教道理呢，你反招了我们来大玩 大笑的！”}
}
{
    \eA{and smoothed it down with a few touches. Afterwards, laying the brush in its
        place she stepped into the outer suite. "Is this," she said pointing at Li
        Wan, "doing what you're told and showing us how to do needlework and
        teaching us manners? Why, instead of that, you press us to come here and
        have a good romp and a hearty laugh!" "Just you listen to her perverse
        talk," Li Wan laughed. "She takes the lead and kicks up a rumpus,}
}
{
}

\Verse{200}
{
    \zA{李纨笑道：“你们听他这刁话。}
}
{
    \eA{and incites people to laugh, and then she throws the blame upon me! In real
        truth, she's a despicable thing!}
}
{
}

\Verse{201}
{
    \zA{他领着头儿闹，引着人笑了，倒赖我的不 是!真真恨的我!只保佑你明儿得一个利害婆婆，再得几个千刁万恶的大姑子、小姑 子，试试你那会子还这么刁不刁了！”}
}
{
    \eA{What I wish is that you should soon get some dreadful mother-in-law, and
        several crotchety and abominable older and younger sisters-in-law, and we'll
        see then whether you'll still be as perverse or not!" Tai-yü at once became
        quite scarlet in the face, and pulling Pao-ch'ai, "Let us," she added, "give
        her a whole year's leave!" "I've got an impartial remark to make.}
}
{
}

\Verse{202}
{
    \zA{黛玉早红了脸，拉着宝钗说：“咱们放他一年的假罢。”}
}
{
    \eA{Listen to me all of you!" Pao-ch'ai chimed in. "Albeit the girl, Ou, may
        have some idea about painting, all she can manage are just a few outline
        sketches,}
}
{
}

\Verse{203}
{
    \zA{宝钗道：“我有一句公 道话，你们听听：藕丫头虽会画，不过是几笔写意；如今画这园子，非离了肚子里 头有些丘壑的，如何成画？}
}
{
    \eA{so that unless, now that she has to accomplish the picture of this garden,
        she can lay a claim to some ingenuity, will she ever be able to succeed in
        effecting a painting? This garden resembles a regular picture. The rockeries
        and trees, towers and pavilions, halls and houses are, as far as distances
        and density go,}
}
{
}

\Verse{204}
{
    \zA{这园子却是像画儿一般，山石树木，楼阁房屋，远近疏 密，也不多，也不少，恰恰的是这样。}
}
{
    \eA{neither too numerous, nor too few. Such as it is, it is fitly laid out; but
        were you to put it on paper in strict compliance with the original, why, it
        will surely not elicit admiration. In a thing like this, it's necessary to
        pay due care to the various positions and distances on paper,}
}
{
}

\Verse{205}
{
    \zA{你若照样儿往纸上一画，是必不能讨好的。}
}
{
    \eA{whether they should be large or whether small; and to discriminate between
        main and secondary; adding what is needful to add,}
}
{
}

\Verse{206}
{
    \zA{这要看纸的地步远近，该多该少，分主分宾，该添的要添，该藏该减的要藏要减， 该露的要露，这一起了稿子，再端详斟酌，方成一幅图样。}
}
{
    \eA{concealing and reducing what should be concealed and reduced, and exposing
        to view what should remain visible. As soon as a rough copy is executed, it
        should again be considered in all its details, for then alone will it assume
        the semblance of a picture. In the second place, all these towers, terraces
        and structures must be distinctly delineated;}
}
{
}

\Verse{207}
{
    \zA{第二件：这些楼台房舍， 是必要界划的。}
}
{
    \eA{for with just a trifle of inattention, the railings will slant, the pillars
        will be topsy-turvy, doors and windows will recline in a horizontal
        position,}
}
{
}

\Verse{208}
{
    \zA{一点儿不留神，栏杆也歪了，柱子也塌了，门窗也倒竖过来，阶砌 也离了缝，甚至桌子挤到墙里头去，花盆放在帘子上来，岂不倒成了一张笑话儿了!
        第三：要安插人物，也要有疏密，有高低。}
}
{
    \eA{steps will separate, leaving clefts between them, and even tables will be
        crowded into the walls, and flower-pots piled on portières; and won't it,
        instead of turning out into a picture, be a mere caricature? Thirdly, proper
        care must also be devoted, in the insertion of human beings, to density and
        height, to the creases of clothing,}
}
{
}

\Verse{209}
{
    \zA{衣褶裙带，指手足步，最是要紧；一笔 不细，不是肿了手，就是瘸了脚，染脸撕发倒是小事。}
}
{
    \eA{to jupes and sashes, to fingers, hands, and feet, as these are most
        important details; for if even one stroke be not thoroughly executed, then,
        if the hands be not swollen, the feet will be made to look as if they were
        lame. The colouring of faces and the drawing of the hair are minor points;}
}
{
}

\Verse{210}
{
    \zA{依我看来，竟难的很。}
}
{
    \eA{but, in my own estimation, they really involve intense difficulty. Now a
        year's leave is, on one hand,}
}
{
}

\Verse{211}
{
    \zA{如今 一年的假也太多，一月的假也太少，竟给他半年的假；再派了宝兄弟帮着他。}
}
{
    \eA{too excessive, and a month's is, on the other, too little; so just give her
        half a year's leave. Depute, besides, cousin Pao-yü to lend her a hand in
        her task. Not that cousin Pao knows how to give any hints about painting;}
}
{
}

\Verse{212}
{
    \zA{并不 是为宝兄弟知道教着他画，――那就更误了事；为的是有不知道的，或难安插的， 宝兄弟拿出去问问那会画的先生们，就容易了。”}
}
{
    \eA{that in itself would be more of a drawback; but in order that, in the event
        of there being anything that she doesn't comprehend, or of anything
        perplexing her as to how best to insert it, cousin Pao may take the picture
        outside and make the necessary inquiries of those gentlemen, who excel in
        painting. Matters will thus be facilitated for her."}
}
{
}

\Verse{213}
{
    \zA{宝玉听了，先喜的说：“这话极是。}
}
{
    \eA{At this suggestion Pao-yü was the first to feel quite enchanted. "This
        proposal is first-rate!"}
}
{
}

\Verse{214}
{
    \zA{詹子亮的工细楼台就极好，程日兴的美人是绝技，如今就问他们去。”}
}
{
    \eA{he exclaimed. "The towers and terraces minutely executed by Chan Tzu-liang
        are so perfect, and the beauties painted by Ch'eng Jih-hsing so extremely
        fine that I'll go at once and ask them of them!" "I've always said that you
        fuss for nothing!" Pao-ch'ai interposed. "I merely passed a cursory remark,}
}
{
}

\Verse{215}
{
    \zA{宝钗道：“我说你是‘无事忙’，说了一声，你就问他去!也等着商议定了再去。}
}
{
    \eA{and there you want to go immediately and ask for things. Do wait until we
        arrive at some decision in our deliberations, and then you can go! But let's
        consider now what would be best to use to paint the picture on?"}
}
{
}

\Verse{216}
{
    \zA{如今且说拿什么画？”}
}
{
    \eA{"I've got, in my quarters," Pao-yü answered, "some snow-white,}
}
{
}

\Verse{217}
{
    \zA{宝玉道：“家里有雪浪纸，又大，又托墨。”}
}
{
    \eA{wavy paper, which is both large in size, and proof against ink as well."
        Pao-ch'ai gave a sarcastic smile.}
}
{
}

\Verse{218}
{
    \zA{宝钗冷笑道：“我 说你不中用。}
}
{
    \eA{"I do maintain," she cried, "that you are a perfectly useless creature!}
}
{
}

\Verse{219}
{
    \zA{那雪浪纸写字、画写意画儿，或是会山水的画南宗山水，托墨，禁得 皴染；拿了画这个，又不托色，又难烘，画也不好，纸也可惜。}
}
{
    \eA{That snow-white, wavy paper is good for pictures consisting of characters
        and for outline drawings. Or else, those who have the knack of making
        landscapes, use it for depicting scenery of the southern Sung era, as it
        resists ink and is strong enough to bear coarse painting. But were you to
        employ this sort of paper to make a picture of this garden on,}
}
{
}

\Verse{220}
{
    \zA{我教给你一个法子： 原先盖这园子就有一张细致图样，虽是画工描的，那地步方向是不错的。}
}
{
    \eA{it will neither stand the colours, nor will it be easy to dry the painting
        by the fire. So not only won't it be suitable, but it will be a pity too to
        waste the paper. I'll tell you a way how to get out of this. When this
        garden was first laid out,}
}
{
}

\Verse{221}
{
    \zA{你和太太 要出来，也比着那纸的大小，和凤姐姐要一块重绢，交给外边相公们，叫他照着这 图样删补着立了稿子，添了人物，就是了。}
}
{
    \eA{some detailed plan was used, which although executed by a mere
        house-decorator, was perfect with regard to sites and bearings. You'd better
        therefore ask for it of your worthy mother, and apply as well to lady Feng
        for a piece of thick glazed lustring of the size of that paper, and hand
        them to the gentlemen outside,}
}
{
}

\Verse{222}
{
    \zA{就是配这些青绿颜色，并泥金泥银，也 得他们配去。}
}
{
    \eA{and request them to prepare a rough copy for you, with any alterations or
        additions as might be necessary to make so as to accord with the style of
        these grounds. All that will remain to be done will be to introduce a few
        human beings; no more. Then when you have to match the azure and green
        pigments as well as the ground gold and ground silver,}
}
{
}

\Verse{223}
{
    \zA{你们也得另拢上风炉子，预备化胶、出胶、洗笔。}
}
{
    \eA{you can get those people again to do so for you. But you'll also have to
        bring an extra portable stove, so as to have it handy for melting the glue,}
}
{
}

\Verse{224}
{
    \zA{还得一个粉油大案， 铺上毡子。}
}
{
    \eA{and for washing your pencils, after you've taken the glue off. You further
        require a large table,}
}
{
}

\Verse{225}
{
    \zA{你们那些碟子也不全，笔也不全，都从新再弄一分儿才好。”}
}
{
    \eA{painted white and covered with a cloth. That lot of small dishes you have
        aren't sufficient; your pencils too are not enough. It will be well
        consequently for you to purchase a new set of each."}
}
{
}

\Verse{226}
{
    \zA{惜春道：“我 何曾有这些画器？}
}
{
    \eA{"Do I own such a lot of painting materials!" Hsi Ch'un exclaimed. "Why, I
        simply use any pencil that first comes under my hand to paint with;}
}
{
}

\Verse{227}
{
    \zA{不过随手的笔画画罢了。}
}
{
    \eA{that's all. And as for pigments, I've only got four kinds, ochrey stone,
        'Kuang' flower paint,}
}
{
}

\Verse{228}
{
    \zA{就是颜色，只有赭石、广花、藤黄、胭 脂这四样。}
}
{
    \eA{rattan yellow and rouge. Besides these, all I have amount to a couple of
        pencils for applying colours; no more." "Why didn't you say so earlier?"}
}
{
}

\Verse{229}
{
    \zA{再有不过是两支着色的笔就完了。”}
}
{
    \eA{Pao-ch'ai remarked. "I've still got some of these things remaining. But you
        don't need them, so were I to give you any,}
}
{
}

\Verse{230}
{
    \zA{宝钗道：“你何不早说？}
}
{
    \eA{they'd lie uselessly about. I'll put them away for you now for a time,}
}
{
}

\Verse{231}
{
    \zA{这些东西我 却还有，只是你用不着，给你也白放着。}
}
{
    \eA{and, when you want them, I'll let you have some. You should, however, keep
        them for the exclusive purpose of painting fans; for were you to paint such
        big things with them it would be a pity!}
}
{
}

\Verse{232}
{
    \zA{如今我且替你收着，等你用着这个的时候 我送你些。}
}
{
    \eA{I'll draw out a list for you to-day to enable you to go and apply to our
        worthy senior for the items; as it isn't likely that you people can possibly
        know all that's required. I'll dictate them, and cousin Pao can write them
        down!" Pao-yü had already got a pencil and inkslab ready, for, fearing lest
        he might not remember clearly the various necessaries,}
}
{
}

\Verse{233}
{
    \zA{也只可留着画扇子，若画这大幅的，也就可惜了。}
}
{
    \eA{he had made up his mind to write a memorandum of them; so the moment he
        heard Pao-ch'ai's suggestion, he cheerfully took up his pencil,}
}
{
}

\Verse{234}
{
    \zA{今儿替你开个单子， 照着单子和老太太要去。}
}
{
    \eA{and listened quietly. "Four pencils of the largest size," Pao-ch'ai
        commenced, "four of the third size; four of the second size;}
}
{
}

\Verse{235}
{
    \zA{你们也未必知道的全，我说着，宝兄弟写。”}
}
{
    \eA{four pencils for applying colours on big ground; four on medium ground; four
        for small ground; ten claws of large southern crabs; ten claws of small
        crabs;}
}
{
}

\Verse{236}
{
    \zA{宝玉早已预备下笔砚了，原怕记不清白，要写了记着，听宝钗如此说，喜的提 起笔来静听。}
}
{
    \eA{ten pencils for painting side-hair and eyebrows; twenty for laying heavy
        colours; twenty for light colours; ten for painting faces; twenty
        willow-twigs; four ounces of 'arrow head' pearls; four ounces of southern
        ochre;}
}
{
}

\Verse{237}
{
    \zA{宝钗说道：“头号排笔四支，二号排笔四支，三号排笔四支，大染四 支，中染四支，小染四支，大南蟹爪十支，小蟹爪十支，须眉十支，大着色二十支，
        小着色二十支，开面十支，柳条二十支，箭头朱四两，南赭四面，石黄四两，石青 四两，石绿四两，管黄四两，广花八两，铅粉十四匣，胭脂十二帖，大赤二百帖，
        青金二百帖，广匀胶四两，净矾四两。}
}
{
    \eA{four ounces of stone yellow; four ounces of dark green; four ounces of
        malachite; four ounces of tube-yellow; eight ounces of 'kuang' flower; four
        boxes of lead powder; ten sheets of rouge; two hundred sheets of thin
        red-gold leaves;}
}
{
}

\Verse{238}
{
    \zA{矾绢的胶矾在外，别管他们，只把绢交出去， 叫他们矾去。}
}
{
    \eA{two hundred sheets of lead; four ounces of smooth glue, from the two Kuang;
        and four ounces of pure alum. The glue and alum for sizing the lustring are
        not included,}
}
{
}

\Verse{239}
{
    \zA{这些颜色，咱们淘澄飞跌着，又玩了，又使了，包你一辈子都够使了。}
}
{
    \eA{so don't bother yourselves about them, but just take the lustring and give
        it to them outside to size it with alum for you. You and I can scour and
        clarify all these pigments, and thus amuse ourselves,}
}
{
}

\Verse{240}
{
    \zA{再要顶细绢箩四个，粗箩二个，担笔四支，大小乳钵四个，大粗碗二十个，五寸碟 子十个，三寸粗白碟子二十个，风炉两个，沙锅大小四个，新磁缸二口，新水桶二
        只，一尺长白布口袋四个，浮炭二十斤，柳木炭一二斤，三屉木箱一个，实地纱一 丈，生姜二两，酱半斤――”黛玉忙笑道：“铁锅一口，铁铲一个。”}
}
{
    \eA{and prepare them for use as well. I feel sure you'll have an ample supply to
        last you a whole lifetime. But you must also get ready four sieves of fine
        lustring; a pair of coarse ones; four brush-pencils; four bowls, some large,
        some small; twenty large, coarse saucers; ten five-inch plates; twenty
        three-inch coarse,}
}
{
}

\Verse{241}
{
    \zA{宝钗道：“这 做什么？”}
}
{
    \eA{white plates; two stoves; four large and small earthenware pans;}
}
{
}

\Verse{242}
{
    \zA{黛玉道：“你要生姜和酱这些作料，我替你要铁锅来，好炒颜色吃啊。”}
}
{
    \eA{two new porcelain jars; four new water buckets; four one-foot-long bags,
        made of white cloth; two catties of light charcoal; one or two catties of
        willow-wood charcoal; a wooden box with three drawers;}
}
{
}

\Verse{243}
{
    \zA{众人都笑起来。}
}
{
    \eA{a yard of thick gauze, two ounces of fresh ginger;}
}
{
}

\Verse{244}
{
    \zA{宝钗笑道：“颦儿你知道什么!那粗磁碟子保不住不上火烤，不拿姜 汁子和酱预先抹在底子上烤过，一经了火，是要炸的。”}
}
{
    \eA{half a catty of soy;…" "An iron kettle and an iron shovel," hastily chimed
        in Tai-yü with a smile full of irony. "To do what with them?" Pao-ch'ai
        inquired. "You ask for fresh ginger, soy and all these condiments, so I
        indent for an iron kettle for you to cook the paints and eat them." Tai-yü
        answered, to the intense merriment of one and all,}
}
{
}

\Verse{245}
{
    \zA{众人听说，都道：“这就是 了。”}
}
{
    \eA{who gave way to laughter. "What do you, P'in Erh, know about these things?"}
}
{
}

\Verse{246}
{
    \zA{黛玉又看了一回单子，笑着拉探春悄悄的道：“你瞧瞧，画个画儿，又要起这 些水缸箱子来。}
}
{
    \eA{Pao-ch'ai laughed. "I am not certain in my mind that you won't put those
        coarse coloured plates straightway on the fire. But unless you take the
        precaution beforehand of rubbing the bottom with ginger juice, mixed with
        soy, and of warming them dry, they're bound to crack, the moment they
        experience the least heat."}
}
{
}

\Verse{247}
{
    \zA{想必糊涂了，把他的嫁妆单子也写上了。”}
}
{
    \eA{"It's really so," they exclaimed with one voice, after this explanation.
        Tai-yü perused the list for a while. She then smiled and gave T'an Ch'un a
        tug.}
}
{
}

\Verse{248}
{
    \zA{探春听了，笑个不住， 说道“宝姐姐，你还不拧他的嘴？}
}
{
    \eA{"Just see," she whispered, "we want to paint a picture, and she goes on
        indenting for a number of water jars and boxes! But, I presume, she's got so
        muddled,}
}
{
}

\Verse{249}
{
    \zA{你问问他编派你的话！”}
}
{
    \eA{that she inserts a list of articles needed for her trousseau." T'an Ch'un,
        at her remark,}
}
{
}

\Verse{250}
{
    \zA{宝钗笑道：“不用问，‘狗 嘴里还有象牙不成’！”}
}
{
    \eA{laughed with such heartiness, that it was all she could do to check herself.
        "Cousin Pao," she observed, "don't you wring her mouth?}
}
{
}

\Verse{251}
{
    \zA{一面说，一面走上来，把黛玉按在炕上，便要拧他的脸。}
}
{
    \eA{Just ask her what disparaging things she said about you." "Why need I ask?"
        Pao-ch'ai smiled. "Is it likely, pray, that you can get ivory out of a cur's
        mouth?" Speaking the while, she drew near,}
}
{
}

\Verse{252}
{
    \zA{黛 玉笑着，忙央告道：“好姐姐!饶了我罢!颦儿年纪小，只知说，不知道轻重，做姐 姐的教导我。}
}
{
    \eA{and, seizing Tai-yü, she pressed her down on the stove-couch with the
        intention of pinching her face. Tai-yü smilingly hastened to implore for
        grace. "My dear cousin," she cried, "spare me! P'in Erh is young in years;
        all she knows is to talk at random; she has no idea of what's proper and
        what's improper.}
}
{
}

\Verse{253}
{
    \zA{姐姐不饶我，我还求谁去呢？”}
}
{
    \eA{But you are my elder cousin, so teach me how to behave. If you, cousin,
        don't let me off,}
}
{
}

\Verse{254}
{
    \zA{众人不知话内有因，都笑道：“说的 好可怜见儿的!连我们也软了，饶了他罢。”}
}
{
    \eA{to whom can I go and address my entreaties?" Little did, however, all who
        heard her apprehend that there lurked some hidden purpose in her
        insinuations. "She's right there," they consequently pleaded smilingly.}
}
{
}

\Verse{255}
{
    \zA{宝钗原是和他玩，忽听他又拉扯上前番 说他胡看杂书的话，便不好再和他闹了，放起他来。}
}
{
    \eA{"So much is she to be pitied that even we have been mollified; do spare her
        and finish!" Pao-ch'ai had, at first, meant to play with her, but when she
        unawares heard her drag in again the advice she had tendered her the other
        day, with regard to the reckless perusal of unwholesome books,}
}
{
}

\Verse{256}
{
    \zA{黛玉笑道：“到底是姐姐，要 是我，再不饶人的。”}
}
{
    \eA{she at once felt as if she could not have any farther fuss with her, and she
        let her rise to her feet. "It's you, after all, elder cousin,"}
}
{
}

\Verse{257}
{
    \zA{宝钗笑指他道：“怪不得老太太疼你，众人爱你，今儿我也怪 疼你的了。}
}
{
    \eA{Tai-yü laughed. "Had it been I, I wouldn't have let any one off." Pao-ch'ai
        smiled and pointed at her. "It is no wonder," she said, "that our dear
        ancestor doats on you and that every one loves you.}
}
{
}

\Verse{258}
{
    \zA{过来，我替你把头发笼笼罢。”}
}
{
    \eA{Even I have to-day felt my heart warm towards you! But come here and let me
        put your hair up for you!"}
}
{
}

\Verse{259}
{
    \zA{黛玉果然转过身来，宝钗用手笼上去。}
}
{
    \eA{Tai-yü then, in very deed, swung herself round and crossed over to her.
        Pao-ch'ai arranged her coiffure with her hands.}
}
{
}

\Verse{260}
{
    \zA{宝玉在旁看着，只觉更好，不觉后悔：“不该令他抿上鬓去，也该留着，此时叫他 替他抿上去。”}
}
{
    \eA{Pao-yü, who stood by and looked on, thought the style, in which her hair was
        being made up, better than it was before. But, of a sudden, he felt sorry at
        what had happened, as he fancied that she should not have let her brush her
        side hair,}
}
{
}

\Verse{261}
{
    \zA{正自胡想，只见宝钗说道：“写完了，明儿回老太太去。}
}
{
    \eA{but left it alone for the time being and asked him to do it for her. While,
        however, he gave way to these erratic thoughts, he heard Pao-ch'ai speak.}
}
{
}

\Verse{262}
{
    \zA{若家里有的 就罢，若没有的，就拿些钱去买了来，我帮着你们配。”}
}
{
    \eA{"We've done with what there was to write," she said, "so you'd better
        tomorrow go and tell grandmother about the things. If there be any at home,}
}
{
}

\Verse{263}
{
    \zA{宝玉忙收了单子。}
}
{
    \eA{well and good; but if not, get some money to buy them with.}
}
{
}

\Verse{264}
{
    \zA{大家又说了一回闲话儿。}
}
{
    \eA{I'll then help you both in your preparations." Pao-yü vehemently put the
        list away;}
}
{
}

\Verse{265}
{
    \zA{至晚饭后，又往贾母处来请安。}
}
{
    \eA{after which, they all joined in a further chat on irrelevant matters; and,
        their evening meal over,}
}
{
}

\Verse{266}
{
    \zA{贾母原没有大病，不 过是劳乏了，兼着了些凉，温存了一日，又吃了一两剂药，发散了发散，至晚也就 好了。}
}
{
    \eA{they once more repaired into old lady Chia's apartments to wish her
        good-night. Their grandmother had, indeed, had nothing serious the matter
        with her. Her ailment had amounted mainly to fatigue, to which a slight
        chill had been super-added, so that having kept in the warm room for the day
        and taken a dose or two of medicine,}
}
{
}

\Verse{267}
{
    \zA{不知次日又有何话，下回分解。}
}
{
    \eA{she entirely got over the effects, and felt, in the evening, quite like own
        self again. But, reader, the occurrences of the next day areas yet a mystery
        to you,}
}
{
}

\Verse{268}
{
    \zA{(本章完)}
}
{
    \eA{but the nest chapter will divulge them.}
}
{
}
